% === Macro definitions from preamble ===
\newcommand{\starlanguage}{Significance Indicator: $^{***}
\newcommand{\aisub}{AI agent}
\newcommand{\persona}{prompt}
\newcommand{\Persona}{Prompt}
\newcommand{\NGamesS}{1,500}
\newcommand{\NParticipantS}{4,500}
\newcommand{\TotalGames}{883,320}
\newcommand{\green}[1]{\textcolor{ForestGreen}
\newcommand{\red}[1]{\textcolor{Red}
\newcommand{\WorkingTitle}{General Social Agents}
% === End macro definitions ===

  
\date{\today}
\title{\WorkingTitle{}\thanks{
Thanks to Tyler Cowen and the Mercatus Center for generous funding and intellectual support.
Thanks to 
Alex Moehring,
Daniel Rock,
David Holtz,
Drew Fudenberg,
Jenny Allen,
Jessica Hullman,  
John List,
Kehang Zhu,
Leland Bybee,
Michael Zhao,
Seth Benzell,
Sophia Kazinnik,
Soumitra Shukla, and
Steve Tadelis
for their time and helpful comments.
We are deeply grateful to Reanna Ishmael for software development support.
The experiments in this paper were preregistered on \url{https://aspredicted.org/} numbers $222695$, $231091$, and $241394$.
Author contact information, code, and data are currently or will be available at \url{http://www.benjaminmanning.io/}. 
Both authors have a financial interest in \url{https://www.expectedparrot.com/}. Horton is an economic advisor to Anthropic.
In preparing this paper, the authors utilized generative AI models extensively as tools to assist with editing and evaluation. 
The authors retain full responsibility for all content and conclusions presented herein.
}
} 
\author{
Benjamin S. Manning \\ MIT  \and   
John J. Horton \\ MIT  \& NBER
}
% \date{\today\\\vspace{2em}MOST RECENT DRAFT \href{https://benjaminmanning.io/files/optimize.pdf}{[HERE]}}
\date{\today}

\pagenumbering{gobble} 
\maketitle

\vspace{-0.5cm}
\begin{abstract}

\noindent Useful social science theories predict behavior across settings. 
However, applying a theory to make predictions in new settings is challenging: rarely can it be done without ad hoc modifications to account for setting-specific factors.
We argue that AI agents put in simulations of those novel settings offer an alternative for applying theory, requiring minimal or no modifications.
We present an approach for building such ``general'' agents that use theory-grounded natural language instructions, existing empirical data, and knowledge acquired by the underlying AI during training.
To demonstrate the approach in settings where no data from that data-generating process exists---as is often the case in applied prediction problems---we design a heterogeneous population of \TotalGames{} novel games.
AI agents are constructed using human data from a small set of conceptually related but structurally distinct ``seed'' games.
In preregistered experiments, on average, agents predict initial human play in a random sample of \NGamesS{} games from the population better than (i) a cognitive hierarchy model, (ii) game-theoretic equilibria, and (iii) out-of-the-box agents.
For a small set of separate novel games, these simulations predict responses from a new sample of human subjects \emph{better} even than the most plausibly relevant published human data.
\end{abstract}

\onehalfspacing
\newpage \clearpage
\pagenumbering{arabic}
